\documentclass[11pt,a4paper]{article}

% ── Codificación, fuentes y lenguaje ───────────────────────────────────────────
\usepackage[utf8]{inputenc}
\usepackage[T1]{fontenc}
\usepackage[default,scale=0.95]{sourcesanspro}
\renewcommand{\familydefault}{\sfdefault}
\usepackage[spanish,es-noshorthands,es-tabla]{babel}

% ── Geometría y espaciado ─────────────────────────────────────────────────────
\usepackage[top=2.2cm, bottom=2.5cm, left=2.2cm, right=2.2cm, headheight=15pt]{geometry}
\usepackage{setspace}
\setstretch{1.18}
\usepackage{parskip}

% ── Matemáticas y unidades ────────────────────────────────────────────────────
\usepackage{amsmath}
\usepackage{amssymb}
\usepackage{siunitx}

% ── Circuitos ─────────────────────────────────────────────────────────────────
\usepackage[european, straightvoltages]{circuitikz}

% ── Tablas ────────────────────────────────────────────────────────────────────
\usepackage{booktabs}
\usepackage{tabularx}
\usepackage{array}
\usepackage{longtable}
\usepackage{colortbl}
\usepackage{enumitem}

% ── Colores, cajas e iconos ───────────────────────────────────────────────────
\usepackage[dvipsnames]{xcolor}
\usepackage{tcolorbox}
\tcbuselibrary{skins, breakable}
\usepackage{fontawesome5}
\usepackage{graphicx}

% ── Listados de código (dark-mode infográfico) ────────────────────────────────
%   Estilo base: fondo oscuro con números de línea. Los bloques lstlisting se
%   envuelven automáticamente en una caja tcolorbox redondeada.
\usepackage{listings}

% ── Secciones con barra de color (infografía) ────────────────────────────────
\usepackage{titlesec}
\titleformat{\section}
  {\normalfont\LARGE\bfseries\color{azulNoche}}
  {\colorbox{cyanNeon}{\color{fondoOscuro}\thesection}}{0.7em}{}
  [\vspace{2pt}\color{cyanNeon}\rule{\linewidth}{1.8pt}]
\titlespacing*{\section}{0pt}{24pt}{10pt}
\titleformat{\subsection}
  {\normalfont\large\bfseries\color{azulNoche}}
  {\textcolor{cyanNeon}{\thesubsection}}{0.7em}{}
  [\vspace{1pt}\color{grisLinea}\rule{\linewidth}{0.6pt}]
\titlespacing*{\subsection}{0pt}{14pt}{6pt}
\titleformat{\subsubsection}
  {\normalfont\normalsize\bfseries\color{azulNoche}}
  {\textcolor{cyanNeon}{\thesubsubsection}}{0.7em}{}
\titlespacing*{\subsubsection}{0pt}{10pt}{4pt}

% ── Cabeceras y pies de página ────────────────────────────────────────────────
\usepackage{fancyhdr}
\usepackage{lastpage}
\pagestyle{fancy}
\fancyhf{}
\renewcommand{\headrulewidth}{0pt}
\renewcommand{\footrulewidth}{0pt}
\fancyfoot[C]{%
  \begin{minipage}{\linewidth}
    \vspace{2pt}
    \color{cyanNeon}\rule{\linewidth}{1.2pt}\\[2pt]
    \centering\color{grisTexto}\footnotesize
    Página \thepage\ de \pageref{LastPage}
  \end{minipage}%
}

% ── Hiperenlaces ──────────────────────────────────────────────────────────────
\usepackage[colorlinks=true, linkcolor=azulNoche, urlcolor=cyanNeon]{hyperref}
\sloppy
\emergencystretch 3em

% ── Paleta de colores — tecnológico dark-mode ──────────────────────────────────
\definecolor{fondoOscuro}{HTML}{0D1B2A}     % Dark navy — fondo banda título
\definecolor{verdeTurquesa}{HTML}{004D40}    % Verde turquesa — bandas de marca
\definecolor{azulNoche}{HTML}{1B3A6B}       % Midnight blue — secciones, marcos
\definecolor{azulMedio}{HTML}{2A5298}       % Azul medio — variante de acento
\definecolor{cyanNeon}{HTML}{00C8E0}        % Cyan eléctrico — acentos primarios
\definecolor{verdeSignal}{HTML}{00B887}     % Verde señal — fortalezas, éxito
\definecolor{naranjaVivo}{HTML}{FF7043}     % Naranja vivo — mejoras, advertencias
\definecolor{rojoAlerta}{HTML}{E53935}      % Rojo alerta — errores
\definecolor{grisPapel}{HTML}{F4F6FA}       % Fondo tarjetas claras
\definecolor{grisLinea}{HTML}{C8D0E0}       % Bordes sutiles
\definecolor{grisTexto}{HTML}{6B7A99}       % Texto secundario
\definecolor{amarilloNota}{HTML}{FFB300}    % Notas / advertencias doradas

% Colores específicos para bloques de código (dark-mode)
\definecolor{codigoFondo}{HTML}{1E2D3D}      % Fondo del bloque de código
\definecolor{codigoComentario}{HTML}{7A8FA6} % Comentarios
\definecolor{codigoCadena}{HTML}{FF9D5C}     % Cadenas / strings
\definecolor{codigoKeyword}{HTML}{00C8E0}    % Palabras clave
\definecolor{codigoNumero}{HTML}{00E5A0}     % Números

% ── Banda de título: portada infográfica ──────────────────────────────────────
%   Diseño en 2 capas: banda de color (por defecto fondo oscuro) + franja de
%   acento cyan abajo. Uso: \bandaTitulo[color]{título}{subtítulo}.
%   El icono se puede cambiar con \renewcommand{\iconoBanda}{...} (FontAwesome).
\newcommand{\iconoBanda}{microchip}
\newcommand{\bandaTitulo}[3][fondoOscuro]{%
  \begin{tcolorbox}[
    enhanced,
    colback=#1, colframe=#1,
    boxrule=0pt, arc=8pt, outer arc=8pt,
    width=\linewidth,
    top=18pt, bottom=6pt, left=14pt, right=14pt,
    borderline south={3.5pt}{0pt}{cyanNeon},
  ]
    \begin{minipage}[c]{0.07\linewidth}
      \centering
      \textcolor{cyanNeon}{\fontsize{30}{30}\selectfont\faIcon{\iconoBanda}}
    \end{minipage}%
    \hspace{8pt}%
    \begin{minipage}[c]{0.88\linewidth}
      {\fontsize{20}{24}\selectfont\bfseries\color{white}#2}\par\vspace{5pt}%
      \textcolor{cyanNeon!80!white}{\small\faIcon{angle-right}~#3}%
    \end{minipage}
    \vspace{4pt}
  \end{tcolorbox}%
  \vspace{8pt}%
}

% ── Tarjetas de datos (infografía) ────────────────────────────────────────────
\tcbset{
  tarjetaDato/.style={
    enhanced, breakable,
    arc=6pt, outer arc=6pt,
    colback=grisPapel, colframe=azulNoche,
    boxrule=1pt,
    fonttitle=\bfseries\small, coltitle=white,
    attach boxed title to top left={yshift=-3mm, xshift=6mm},
    boxed title style={
      colback=azulNoche, colframe=azulNoche,
      arc=4pt, boxrule=0pt,
    },
    drop fuzzy shadow=azulNoche!30!white,
    top=8pt, bottom=8pt, left=10pt, right=10pt,
  },
  tarjetaIcono/.style={
    enhanced, breakable,
    arc=6pt, outer arc=6pt,
    colback=white, colframe=grisLinea, boxrule=0.8pt,
    fonttitle=\bfseries\small, coltitle=azulNoche,
    borderline west={3pt}{0pt}{cyanNeon},
    drop fuzzy shadow=azulNoche!25!white,
    top=6pt, bottom=6pt, left=10pt, right=10pt,
  },
  cajaTitulo/.style={
    enhanced, breakable,
    arc=6pt, outer arc=6pt,
    colback=azulNoche!10!grisPapel, colframe=azulNoche,
    boxrule=1pt,
    fonttitle=\bfseries, coltitle=white,
    attach boxed title to top left={yshift=-3mm, xshift=6mm},
    boxed title style={
      colback=azulNoche, colframe=azulNoche,
      arc=4pt, boxrule=0pt,
    },
    drop fuzzy shadow=azulNoche!25!white,
    top=8pt, bottom=8pt, left=10pt, right=10pt,
  },
  cajaContenido/.style={
    enhanced, breakable,
    arc=5pt, outer arc=5pt,
    colback=grisPapel, colframe=grisLinea,
    boxrule=0.8pt,
    fonttitle=\bfseries\small, coltitle=azulNoche,
    top=6pt, bottom=6pt, left=10pt, right=10pt,
  },
  cajaFortaleza/.style={
    enhanced, breakable,
    arc=6pt, outer arc=6pt,
    colback=verdeSignal!8!white, colframe=verdeSignal,
    boxrule=1pt,
    borderline west={4pt}{0pt}{verdeSignal},
    fonttitle=\bfseries\small, coltitle=verdeSignal!70!black,
    drop fuzzy shadow=verdeSignal!20!white,
    top=6pt, bottom=6pt, left=10pt, right=10pt,
  },
  cajaMejora/.style={
    enhanced, breakable,
    arc=6pt, outer arc=6pt,
    colback=naranjaVivo!8!white, colframe=naranjaVivo,
    boxrule=1pt,
    borderline west={4pt}{0pt}{naranjaVivo},
    fonttitle=\bfseries\small, coltitle=naranjaVivo!80!black,
    drop fuzzy shadow=naranjaVivo!20!white,
    top=6pt, bottom=6pt, left=10pt, right=10pt,
  },
  cajaRecomendacion/.style={
    enhanced, breakable,
    arc=6pt, outer arc=6pt,
    colback=cyanNeon!6!white, colframe=cyanNeon!60!azulNoche,
    boxrule=1pt,
    borderline west={4pt}{0pt}{cyanNeon},
    fonttitle=\bfseries\small, coltitle=azulNoche,
    drop fuzzy shadow=azulNoche!20!white,
    top=6pt, bottom=6pt, left=10pt, right=10pt,
  },
}

% ── Ícono + texto (encabezados de sección y elementos) ────────────────────────
\newcommand{\iconotexto}[2]{\textcolor{cyanNeon}{\faIcon{#1}}~#2}

% ── Listados de código: estilo dark + auto-envoltura ─────────────────────────
\lstdefinestyle{estiloCodigo}{
    backgroundcolor=\color{codigoFondo},
    basicstyle=\ttfamily\small\color{white},
    commentstyle=\color{codigoComentario}\itshape,
    stringstyle=\color{codigoCadena},
    keywordstyle=\color{codigoKeyword}\bfseries,
    numberstyle=\tiny\color{codigoComentario},
    numbers=left,
    numbersep=10pt,
    showstringspaces=false,
    breaklines=true,
    breakatwhitespace=false,
    tabsize=4,
    frame=none,
    captionpos=b,
    aboveskip=6pt,
    belowskip=4pt,
    xleftmargin=14pt,
    framexleftmargin=14pt,
    literate=
      {á}{{\'a}}1 {é}{{\'e}}1 {í}{{\'i}}1 {ó}{{\'o}}1 {ú}{{\'u}}1
      {Á}{{\'A}}1 {É}{{\'E}}1 {Í}{{\'I}}1 {Ó}{{\'O}}1 {Ú}{{\'U}}1
      {ñ}{{\~n}}1 {Ñ}{{\~N}}1 {ü}{{\"u}}1 {Ü}{{\"U}}1
      {¿}{{?`}}1 {¡}{{!`}}1
}
\lstdefinestyle{estiloPython}{
    style=estiloCodigo,
    language=Python,
}
\lstset{style=estiloCodigo}
\BeforeBeginEnvironment{lstlisting}{%
  \begin{tcolorbox}[
    enhanced, breakable,
    arc=6pt, outer arc=6pt,
    colback=codigoFondo, colframe=azulNoche,
    boxrule=1pt,
    drop fuzzy shadow=azulNoche!25!white,
    top=2pt, bottom=2pt, left=0pt, right=0pt,
  ]%
}
\AfterEndEnvironment{lstlisting}{\end{tcolorbox}}

% ── Cajas de contenido temáticas ──────────────────────────────────────────────
\newtcolorbox{cajaRecuerda}{
    enhanced, breakable,
    arc=6pt, outer arc=6pt,
    colback=azulNoche!10!grisPapel,
    colframe=azulNoche, boxrule=1pt,
    borderline west={4pt}{0pt}{cyanNeon},
    fonttitle=\bfseries\small\color{white},
    title={\faIcon{info-circle}~Recuerda},
    attach boxed title to top left={yshift=-3mm, xshift=6mm},
    boxed title style={colback=azulNoche, arc=4pt, boxrule=0pt},
    drop fuzzy shadow=azulNoche!25!white,
    left=10pt, right=10pt, top=8pt, bottom=8pt,
}

\newtcolorbox{cajaConcepto}{
    enhanced, breakable,
    arc=6pt, outer arc=6pt,
    colback=verdeSignal!8!white,
    colframe=verdeSignal, boxrule=1pt,
    borderline west={4pt}{0pt}{verdeSignal},
    fonttitle=\bfseries\small\color{white},
    title={\faIcon{lightbulb}~Concepto clave},
    attach boxed title to top left={yshift=-3mm, xshift=6mm},
    boxed title style={colback=verdeSignal!80!black, arc=4pt, boxrule=0pt},
    drop fuzzy shadow=verdeSignal!20!white,
    left=10pt, right=10pt, top=8pt, bottom=8pt,
}

\newtcolorbox{cajaEjemplo}{
    enhanced, breakable,
    arc=6pt, outer arc=6pt,
    colback=cyanNeon!5!white,
    colframe=cyanNeon!70!azulNoche, boxrule=1pt,
    borderline west={4pt}{0pt}{cyanNeon},
    fonttitle=\bfseries\small\color{fondoOscuro},
    title={\faIcon{code}~Ejemplo},
    attach boxed title to top left={yshift=-3mm, xshift=6mm},
    boxed title style={colback=cyanNeon, arc=4pt, boxrule=0pt},
    drop fuzzy shadow=azulNoche!20!white,
    left=10pt, right=10pt, top=8pt, bottom=8pt,
}

\newtcolorbox{cajaEjercicio}{
    enhanced, breakable,
    arc=6pt, outer arc=6pt,
    colback=azulMedio!7!white,
    colframe=azulMedio, boxrule=1pt,
    borderline west={4pt}{0pt}{azulMedio},
    fonttitle=\bfseries\small\color{white},
    title={\faIcon{pencil-alt}~Ejercicio},
    attach boxed title to top left={yshift=-3mm, xshift=6mm},
    boxed title style={colback=azulMedio, arc=4pt, boxrule=0pt},
    drop fuzzy shadow=azulNoche!20!white,
    left=10pt, right=10pt, top=8pt, bottom=8pt,
}

\newtcolorbox{cajaReto}{
    enhanced, breakable,
    arc=6pt, outer arc=6pt,
    colback=naranjaVivo!6!white,
    colframe=naranjaVivo, boxrule=1pt,
    borderline west={4pt}{0pt}{naranjaVivo},
    fonttitle=\bfseries\small\color{white},
    title={\faIcon{fire}~Reto},
    attach boxed title to top left={yshift=-3mm, xshift=6mm},
    boxed title style={colback=naranjaVivo, arc=4pt, boxrule=0pt},
    drop fuzzy shadow=naranjaVivo!20!white,
    left=10pt, right=10pt, top=8pt, bottom=8pt,
}

% ── Indicadores de dificultad ─────────────────────────────────────────────────
\newcommand{\facil}{\textcolor{verdeSignal}{\faIcon{circle}\,\textbf{Fácil}}}
\newcommand{\intermedio}{\textcolor{naranjaVivo}{\faIcon{circle}\,\textbf{Intermedio}}}
\newcommand{\dificil}{\textcolor{rojoAlerta}{\faIcon{circle}\,\textbf{Difícil}}}

% ─────────────────────────────────────────────────────────────────────────────

\begin{document}

\fancyhead[L]{\footnotesize\color{grisTexto}Revisión de apuntes \textbullet\ Temas 7 y 8}
\fancyhead[R]{\footnotesize\color{cyanNeon}\faIcon{microchip}~CIRC\_DISP\_ELECT}

\bandaTitulo{Revisión de Apuntes: Temas 7 y 8}{Circuitos y Dispositivos Electrónicos \textbullet\ Prof. Julima Anato \textbullet\ 19/09/2026}

\section{\iconotexto{clipboard-check}{Alcance y metodología}}

Se revisaron los dos apuntes del curso ubicados en \texttt{docs/CIRC\_DISP\_ELECT/}:

\begin{tabularx}{\linewidth}{@{} >{\raggedright\arraybackslash}X >{\raggedright\arraybackslash}X c l @{}}
\toprule
\textbf{Archivo} & \textbf{Título del documento} & \textbf{Págs.} & \textbf{Creado} \\
\midrule
\texttt{Tema\_7\_Respuesta\_AC.pdf} & Respuesta A.C. del BJT & 10 & 18/08/2026 \\
\texttt{Tema\_8\_Modelo\_Híbrido.pdf} & Análisis del BJT con Modelo Híbrido & 11 & 18/08/2026 \\
\bottomrule
\end{tabularx}

Ambos generados con Microsoft Word para Microsoft 365, autoría de la Prof. Julima Anato.

\textbf{Metodología de revisión:}
\begin{itemize}[leftmargin=1.4em]
  \item Extracción completa de texto con \texttt{pdftotext -layout} (poppler).
  \item Verificación algebraica manual de \emph{todas} las derivaciones: rectas de carga d.c./a.c., condición de MDS, procedimiento de diseño (casos A, B y caso especial), parámetros híbridos y fórmulas de $Z_i$, $Z_o$, $A_v$, $A_i$.
  \item Comprobación de referencias cruzadas (figuras, secciones y ecuaciones citadas en el texto).
  \item Revisión tipográfica y gramatical.
\end{itemize}

\begin{cajaRecuerda}
\texttt{pdftotext} no extrae los glifos de la fuente \emph{Symbol} de Word ($\beta$, $\Omega$, $\approx$, $\leq$, $\ll$); en el texto extraído aparecen como espacios en blanco. Su presencia se confirmó por contexto en los lugares esperados, por lo que \textbf{no} se reportan como errores del documento.
\end{cajaRecuerda}

\section{\iconotexto{wave-square}{Tema 7 --- Respuesta A.C. del BJT}}

\subsection{Contenido}
\begin{itemize}[leftmargin=1.4em]
  \item Amplificador emisor común con $C_E$: función de los condensadores de acoplo ($C_1$, $C_2$) y de desacoplo ($C_E$).
  \item Análisis por superposición: análisis d.c. (polarización por divisor de voltaje, equivalente de Thévenin) y análisis a.c. (circuito equivalente, ecuación de salida).
  \item Rectas de carga d.c. y a.c.; análisis gráfico y ubicación del punto Q.
  \item Máximo Desplazamiento Simétrico (MDS): con $C_E$ y sin $C_E$.
  \item Procedimiento de diseño para MDS: casos A y B, y caso especial sin punto Q conocido.
\end{itemize}

\subsection{Verificación matemática}

\begin{tcolorbox}[cajaFortaleza, title={\faIcon{check-circle}~Derivaciones verificadas --- todas correctas}]
\begin{align*}
\text{Salida a.c. (con } C_E\text{):}\quad
  & v_{CE} = -\,i_C\,(R_C \parallel R_L) \\
\text{Recta de carga a.c.:}\quad
  & V_{CE} = V_{CEQ} - (I_C - I_{CQ})(R_C \parallel R_L) \\
\text{Cortes con los ejes:}\quad
  & V_{CE_{\max}} = V_{CEQ} + I_{CQ}(R_C \parallel R_L), \qquad
  I_{C_{\max}} = \frac{V_{CEQ}}{R_C \parallel R_L} + I_{CQ} \\
\text{Condición de MDS:}\quad
  & V_{CEQ} = V_{CE_{sat}} + I_{CQ}(R_C \parallel R_L) \\
\text{Recta a.c. sin } C_E\text{:}\quad
  & V_{CE} = V_{CEQ} - (I_C - I_{CQ})\bigl[(R_C \parallel R_L) + R_E\bigr] \\
\text{MDS sin } C_E\text{:}\quad
  & V_{CEQ} = V_{CE_{sat}} + I_{CQ}\bigl[(R_C \parallel R_L) + R_E\bigr] \\
\text{Diseño, caso B:}\quad
  & R_B = (\beta + 1)\,\frac{R_E}{10}
\end{align*}

El caso especial de diseño (pág.~10, con $R_E = R_C/10$) fue \textbf{re-derivado desde cero} a partir de la recta de carga d.c. $V_{CC} = V_{CEQ} + 1.1\,I_{CQ}R_C$ y de la condición de MDS; el resultado coincide exactamente con el documento:
\[
R_C = R_L\left[\frac{V_{CC} - 2.1\,V_{CEQ} + 1.1\,V_{CE_{sat}}}{1.1\,(V_{CEQ} - V_{CE_{sat}})}\right],
\qquad
V_{CE_{sat}} < V_{CEQ} < \frac{V_{CC} + 1.1\,V_{CE_{sat}}}{2.1}
\]
\end{tcolorbox}

\subsection{Hallazgos}

\begin{tcolorbox}[cajaMejora, title={\faIcon{exclamation-triangle}~8 erratas detectadas --- 1 con impacto conceptual}]
El hallazgo No.~1 es un error de referencia conceptual; el resto son erratas tipográficas o gramaticales que no alteran el contenido técnico. Los PDF no tienen fuente editable en el repositorio (ni \texttt{.tex} ni \texttt{.docx}): las correcciones deben aplicarse en el archivo Word original.
\end{tcolorbox}

\begin{tabularx}{\linewidth}{@{} c l >{\raggedright\arraybackslash}X >{\raggedright\arraybackslash}X @{}}
\toprule
\textbf{No.} & \textbf{Ubicación} & \textbf{Hallazgo} & \textbf{Corrección propuesta} \\
\midrule
\textcolor{rojoAlerta}{\faIcon{exclamation-circle}}~1 & pág.~8, §5.2 & Se afirma que la sustitución ``permite obtener la ecuación de recta de carga \textbf{d.c.}'', pero la ecuación derivada $V_{CE} = V_{CEQ} - (I_C - I_{CQ})[(R_C \parallel R_L) + R_E]$ es la recta de carga \textbf{a.c.} del amplificador sin $C_E$ (en §3.2 la misma derivación sí se denomina a.c.) & Cambiar ``recta de carga d.c.'' por ``recta de carga a.c.'' \\
\addlinespace[2pt]
2 & pág.~1, §2 & ``de tal manera que se bloquee \textbf{la el} paso de la señal'' (duplicación) & ``se bloquee el paso de la señal'' \\
\addlinespace[2pt]
3 & pág.~2, §2 & ``la resistencia $R_E$ permanece \textbf{en sin} cambios'' (duplicación) & ``permanece sin cambios'' \\
\addlinespace[2pt]
4 & pág.~2, §3 & ``se considera lineal a todo el circuito amplificador \textbf{y aplicar} las propiedades de la linealidad'' & ``\dots\ y se pueden aplicar las propiedades de la linealidad'' \\
\addlinespace[2pt]
5 & pág.~2, §3.1 & ``$C_1$, $C_2$ y $C_E \rightarrow$ \textbf{se actúan} como circuito abierto'' & ``actúan como circuito abierto'' \\
\addlinespace[2pt]
6 & pág.~4, §3.3 & ``$R_C$ // \textbf{Rl}'' con ele minúscula, inconsistente con $R_L$ & $R_C \parallel R_L$ \\
\addlinespace[2pt]
7 & págs.~5--6, §4 & ``está \textbf{mas} cerca'' (dos apariciones, falta la tilde) & ``está más cerca'' \\
\addlinespace[2pt]
8 & pág.~8, §5.2 & ``\textbf{Como resultado de} la superposición \textbf{indica que}'', construcción gramatical mezclada & ``Como resultado de la superposición se tiene que'' (o ``La superposición indica que'') \\
\bottomrule
\end{tabularx}

\section{\iconotexto{project-diagram}{Tema 8 --- Análisis del BJT con Modelo Híbrido}}

\subsection{Contenido}
\begin{itemize}[leftmargin=1.4em]
  \item Modelos circuitales del BJT (modelo $r_e$, híbrido, Ebers--Moll); el híbrido como el más utilizado.
  \item Teoría de cuadripolos: obtención de los parámetros $h_{ie}$, $h_{re}$, $h_{fe}$, $h_{oe}$ y modelo híbrido en emisor común (completo y simplificado).
  \item Amplificador emisor común \textbf{con} $C_E$: cálculo de $Z_i$, $Z_o$, $A_v$, $A_i$.
  \item Amplificador emisor común \textbf{sin} $C_E$: método de ``reflexión hacia la base'' y cálculo de $Z_i$, $Z_o$, $A_v$, $A_i$.
\end{itemize}

\subsection{Verificación matemática}

\begin{tcolorbox}[cajaFortaleza, title={\faIcon{check-circle}~Fórmulas verificadas --- todas correctas}]
Definiciones de los parámetros híbridos (coinciden con el estándar de cuadripolos):
\[
h_{ie} = \left.\frac{v_{BE}}{i_B}\right|_{v_{CE}=0}, \quad
h_{re} = \left.\frac{v_{BE}}{v_{CE}}\right|_{i_B=0}, \quad
h_{fe} = \left.\frac{i_C}{i_B}\right|_{v_{CE}=0}, \quad
h_{oe} = \left.\frac{i_C}{v_{CE}}\right|_{i_B=0}
\]

\textbf{Emisor común con $C_E$} (modelo simplificado, $h_{re} \approx 0$, $h_{oe} \approx 0$):
\[
Z_i = R_B \parallel h_{ie}, \qquad
Z_o = R_C, \qquad
A_v = -\,\frac{h_{fe}\,(R_C \parallel R_L)}{h_{ie}}, \qquad
A_i = A_v\,\frac{Z_i}{R_L}
\]

\textbf{Emisor común sin $C_E$} (reflexión hacia la base, $R_E$ se refleja como $(h_{fe}+1)R_E$):
\[
Z_i = R_B \parallel \bigl[h_{ie} + (h_{fe}+1)R_E\bigr], \qquad
Z_o = R_C
\]
\[
A_v = -\,\frac{h_{fe}\,(R_C \parallel R_L)}{h_{ie} + (h_{fe}+1)R_E}, \qquad
A_i = A_v\,\frac{Z_i}{R_L}
\]

El signo negativo de $A_v$ (desfasaje de $180^{\circ}$ en emisor común) y la anulación de $h_{fe}i_B$ al hacer $v_i = 0$ para el cálculo de $Z_o$ están correctamente justificados en el texto.
\end{tcolorbox}

\subsection{Hallazgos}

\begin{tcolorbox}[cajaMejora, title={\faIcon{exclamation-triangle}~3 erratas detectadas --- 1 de referencia cruzada}]
Igual que en el Tema 7, las correcciones deben aplicarse en el archivo Word original.
\end{tcolorbox}

\begin{tabularx}{\linewidth}{@{} c l >{\raggedright\arraybackslash}X >{\raggedright\arraybackslash}X @{}}
\toprule
\textbf{No.} & \textbf{Ubicación} & \textbf{Hallazgo} & \textbf{Corrección propuesta} \\
\midrule
\textcolor{rojoAlerta}{\faIcon{exclamation-circle}}~1 & pág.~10, §5 (cálculo de $Z_o$) & Dice ``para el cálculo de $Z_o$ en el circuito de la \textbf{figura 10}'', pero el circuito que se está analizando es el de la \textbf{figura 18} (con reflexión hacia la base); la figura~10 corresponde al amplificador \emph{con} $C_E$ del §4 & Cambiar ``figura 10'' por ``figura 18'' \\
\addlinespace[2pt]
2 & pág.~9, §5 & ``un esquema mucho \textbf{mas} sencillo de analizar'' (falta la tilde) & ``mucho más sencillo'' \\
\addlinespace[2pt]
3 & pág.~4, §3 & ``En $h_{oe}$, $i_C \ll v_{CE}$ por lo que $h_{oe}$ resulta en una admitancia cero'' --- redacción confusa (compara magnitudes, no el cociente) & Sugerido: ``la relación $i_C / v_{CE}$ es muy pequeña, por lo que $h_{oe} \approx 0$'' (no es un error formal) \\
\bottomrule
\end{tabularx}

\section{\iconotexto{flag-checkered}{Conclusiones generales}}

\begin{tcolorbox}[cajaRecomendacion, title={\faIcon{check-double}~Veredicto}]
\begin{itemize}[leftmargin=1.4em]
  \item \textbf{Ambos documentos son técnicamente sólidos:} todas las derivaciones matemáticas fueron verificadas (y el caso especial de diseño del Tema 7 re-derivado desde cero) sin discrepancias.
  \item \textbf{Prioridad de corrección} (en el fuente Word de la profesora): los dos errores de referencia --- Tema 7, pág.~8 §5.2 (``recta de carga d.c.'' $\rightarrow$ ``a.c.'') y Tema 8, pág.~10 §5 (``figura 10'' $\rightarrow$ ``figura 18'') --- pues pueden inducir a confusión al estudiante.
  \item Las demás erratas (duplicaciones de palabras, tildes en ``más'', notación $Rl$, gramática) son menores y cosméticas; pueden corregirse en una próxima revisión del material.
\end{itemize}
\end{tcolorbox}

\end{document}
